% 06_tablas/matriz_compatibilidad.tex
% Tabla 4. Matriz de Verificación de Compatibilidad Constitucional y Ambiental
% Capítulo IV, sección 4.2.2 — Propuesta concreta
% Usa longtable para permitir extensión en varias páginas

\begingroup
\small
\begin{longtable}{|p{0.03\textwidth}|p{0.17\textwidth}|p{0.23\textwidth}|p{0.27\textwidth}|p{0.21\textwidth}|}

\caption{Matriz de Verificación de Compatibilidad Constitucional y Ambiental para Proyectos de Ley de Apertura Minera Privada}
\label{tab:compatibilidad} \\

\hline
\textbf{N.°} & \textbf{Criterio de compatibilidad} & \textbf{Fundamento constitucional, legal o jurisprudencial} & \textbf{Estándar exigible mínimo} & \textbf{Indicador de verificación} \\
\hline
\endfirsthead

\multicolumn{5}{l}{\small\textit{(continuación de la Tabla \ref*{tab:compatibilidad})}} \\[2pt]
\hline
\textbf{N.°} & \textbf{Criterio de compatibilidad} & \textbf{Fundamento constitucional, legal o jurisprudencial} & \textbf{Estándar exigible mínimo} & \textbf{Indicador de verificación} \\
\hline
\endhead

\hline
\endfoot

1 &
Derechos de la naturaleza &
Arts.\ 71--74, 395--397, 407 CRE; COA arts.\ 4, 6, 7; Sentencias 1149-19-JP/21 (Bosque Los Cedros) y 22-18-IN/21 (manglares) &
La reforma no suprime mecanismos de protección de la naturaleza como sujeto; contempla la prohibición de actividades en áreas protegidas (art.\ 407 CRE) con excepción calificada; no elimina la obligación de restauración y regeneración de ciclos vitales &
Existe disposición expresa sobre titularidad y exigibilidad de los derechos de la naturaleza. No se derogan disposiciones de protección vigentes. Se respeta la prohibición constitucional del art.\ 407 CRE \\
\hline

2 &
Principio de prevención &
Art.\ 396 CRE; COA arts.\ 9.8, 162, 166, 174; Reglamento COA arts.\ 431--435; Sentencia 1149-19-JP/21; OC-23/17 Corte IDH párr.\ 118--120 &
La reforma mantiene la evaluación ambiental previa como requisito ineludible; no introduce silencio positivo que habilite actividades sin evaluación; no reduce el contenido mínimo del EIA ni el plan de manejo &
No existe disposición de silencio positivo automático en la autorización ambiental. El instrumento de evaluación incluye EIA, plan de manejo y garantías económicas. La Autoridad Ambiental Nacional conserva competencia (COA art.\ 166) \\
\hline

3 &
Principio de precaución &
Arts.\ 73, 395.4 CRE; COA art.\ 9; Sentencia 1149-19-JP/21 (tres elementos: riesgo grave o irreversible, incertidumbre científica, medidas eficaces y oportunas); OC-23/17 Corte IDH párr.\ 130--137 &
La reforma permite la suspensión o restricción de actividades ante incertidumbre científica de daño grave o irreversible; no invierte la carga de la prueba sobre la inocuidad del impacto; mantiene evaluación por subfase del ciclo minero &
Existe mecanismo de suspensión ante riesgo de daño grave e irreversible. La carga de probar la inocuidad del impacto corresponde al operador. No hay transición automática entre subfases sin nueva evaluación ambiental \\
\hline

4 &
Principio de no regresividad ambiental &
Art.\ 161 COA (prohibición expresa); art.\ 11.8 CRE; Sentencia 32-17-IN/21 (reserva de ley orgánica como límite a regulación reglamentaria de derechos) &
La reforma no reduce las obligaciones ambientales sustantivas ni suprime garantías procesales del licenciamiento ambiental; respeta o supera los estándares del COA como ley orgánica ambiental; si simplifica procedimientos, demuestra equivalencia real de los controles alternativos &
Comparación artículo a artículo entre texto anterior y reformado confirma que no se eliminan obligaciones ambientales sustantivas. La simplificación procedimental va acompañada de controles alternativos de eficacia equivalente \\
\hline

5 &
\textit{In dubio pro natura} &
Art.\ 395.4 CRE; COA arts.\ 4, 9; Sentencia 1149-19-JP/21 (ratio decidendi: interpretación en caso de duda a favor de la naturaleza y de las comunidades) &
La reforma incluye o mantiene cláusula de interpretación favorable a la naturaleza; no introduce disposición que invierta el principio; los actos de autorización se interpretan restrictivamente frente a ecosistemas frágiles o vulnerables &
Existe cláusula de interpretación favorable a la naturaleza en disposiciones generales o finales. No hay disposición que prefiera el desarrollo económico sobre la protección ambiental en caso de duda \\
\hline

6 &
Consulta ambiental (art.\ 398 CRE) &
Art.\ 398 CRE; Sentencia 1149-19-JP/21 (nulidad por omisión; estándares: oportunidad, información suficiente, resultado vinculante; consulta realizada por el Estado); Acuerdo de Escazú art.\ 7; COA art.\ 163 &
La reforma prevé la consulta ambiental como acto previo e ineludible a la autorización minera; es realizada por el Estado, no por el concesionario; no la sustituye por declaraciones juradas del titular ni por silencio positivo; cumple los estándares de la Sentencia 1149-19-JP/21 &
La consulta ambiental figura expresamente como acto previo autónomo a la autorización. La autoridad competente es estatal. Se cumplen los estándares de oportunidad, información suficiente y resultado vinculante establecidos por la CCE \\
\hline

7 &
Consulta previa, libre e informada (art.\ 57.7 CRE) &
Art.\ 57.7 CRE; Convenio 169 OIT arts.\ 6, 15; Sentencia 22-18-IN/21 (distinción consulta previa/consulta ambiental; inconstitucionalidad de arts.\ 462--463 RCOAM; reserva de ley orgánica); Acuerdo de Escazú arts.\ 7.2, 7.6 &
La reforma distingue expresamente la CPLI de la consulta ambiental; regula la CPLI por ley orgánica, no por reglamento; activa la CPLI cuando la actividad afecta territorios de comunidades, pueblos y nacionalidades indígenas; no subordina la CPLI al procedimiento de licenciamiento ambiental &
La CPLI está regulada por ley orgánica. Existe mecanismo para identificar si la actividad minera afecta territorios indígenas y activar la CPLI. El cumplimiento de la consulta ambiental no sustituye la CPLI \\
\hline

8 &
Participación ciudadana ambiental &
Arts.\ 61.2, 395.3, 398 CRE; COA arts.\ 18, 163; Acuerdo de Escazú arts.\ 5, 6, 7; Sentencia 22-18-IN/21 (art.\ 184 COAM condicionalmente constitucional con interpretación favorable a participación efectiva) &
La reforma garantiza acceso oportuno a información ambiental; prevé mecanismos de participación en todas las fases del ciclo minero; no reduce los mecanismos participativos del COA; respeta los estándares del Acuerdo de Escazú sobre participación pública en decisiones ambientales &
Los mecanismos de participación ciudadana en el licenciamiento se mantienen. Existe acceso a la información ambiental de los proyectos mineros. No se eliminan los consejos consultivos de la política minera \\
\hline

9 &
Licenciamiento ambiental &
COA arts.\ 162, 166, 174, 179, 181, 183; Reglamento COA arts.\ 431--435; Ley de Minería art.\ 26; Sentencia 1149-19-JP/21 (nulidad de registro ambiental por omisión de consulta en el licenciamiento) &
La reforma mantiene el licenciamiento ambiental (o instrumento equivalente de igual o mayor rigor) como requisito ineludible; no introduce silencio positivo; no reduce el contenido mínimo del instrumento de evaluación; mantiene la competencia de la Autoridad Ambiental Nacional en sectores estratégicos &
El instrumento de evaluación exigido equivale en contenido y rigor al COA art.\ 162. No existe silencio positivo que habilite actividades sin evaluación completa. La Autoridad Ambiental Nacional conserva competencia en sectores estratégicos \\
\hline

10 &
Fiscalización estatal &
Arts.\ 313, 315, 316, 317 CRE; LM arts.\ 8, 9; COA art.\ 166; OC-23/17 Corte IDH párr.\ 142 (obligación estatal de supervisión efectiva de actividades privadas con impacto ambiental) &
La reforma garantiza que el Estado mantenga capacidad efectiva de fiscalización en todas las fases del ciclo minero; no subordina el control ambiental al interés de la entidad promotora de la actividad; establece auditorías periódicas; la capacidad de fiscalización es proporcional a la escala autorizada &
Las atribuciones de control e inspección de los organismos competentes se mantienen. No existe subordinación del órgano de control al ente promotor. Existe dotación institucional suficiente para cubrir el universo de concesiones activas \\
\hline

11 &
Transparencia &
Arts.\ 18, 91, 313 CRE; COA arts.\ 18, 163; Acuerdo de Escazú arts.\ 5, 6; LM art.\ 9 (Catastro y Registro Minero) &
La reforma garantiza acceso público a los EIA y planes de manejo; establece publicación periódica de informes de auditoría; no crea mecanismos de confidencialidad sobre información ambiental relevante para las comunidades; implementa transparencia activa conforme al Acuerdo de Escazú &
Los EIA y las autorizaciones ambientales son accesibles al público. Los informes de auditoría se publican periódicamente. No existen cláusulas de confidencialidad sobre información ambiental de relevancia pública \\
\hline

12 &
Responsabilidad ambiental &
Arts.\ 396, 397 CRE; COA arts.\ 10, 11; LM arts.\ 78--84; COIP arts.\ 257--261; Sentencia 1149-19-JP/21 (parte resolutiva) &
La reforma mantiene la responsabilidad objetiva del concesionario por daños ambientales; exige garantías económicas previas que cubran los costos de remediación; no crea mecanismos de exención que trasladen la responsabilidad al Estado; reconoce la responsabilidad subsidiaria del Estado (art.\ 397 CRE) &
La responsabilidad objetiva del concesionario se mantiene. Las garantías económicas previas son condición de la autorización. No existen cláusulas de exención de responsabilidad ambiental \\
\hline

13 &
Reparación integral &
Arts.\ 397, 11.9 CRE; COA arts.\ 18, 299--325; LM art.\ 27 (cierre de minas); Sentencia 1149-19-JP/21 (parte resolutiva: reparación integral de la naturaleza); Guía CEDEC 2023 (sistematización de criterios de reparación) &
La reforma mantiene la obligación de reparación integral en todas las fases del ciclo minero; contempla el cierre de minas con criterios de reparación integral; las garantías económicas cubren los costos en escenarios de daño grave; no sustituye la restauración por compensaciones monetarias cuando la restauración es posible &
Las dimensiones de la reparación integral (restauración ecológica, rehabilitación social, garantías de no repetición) se mantienen. Las garantías económicas cubren escenarios de daño grave. El plan de manejo incluye monitoreo post-cierre \\
\hline

14 &
Seguridad jurídica &
Art.\ 82 CRE; Sentencia 32-17-IN/21; Sentencia 1149-19-JP/21 (la seguridad jurídica no desplaza la consulta ambiental ni los controles constitucionales) &
La reforma establece normas claras, previas y públicas; define competencias sin solapamientos; prevé procedimientos previsibles de autorización, modificación y extinción de concesiones; distingue la estabilidad del título concesional de los controles ambientales constitucionales, que no son negociables &
Las competencias de MAATE, ARCOM y Ministerio Sectorial están delimitadas. Existe período de transición razonable. La reforma explicita que la seguridad jurídica del concesionario no exime del cumplimiento de las obligaciones ambientales constitucionales \\
\hline

15 &
Diferenciación entre minería formal, informal e ilegal &
Arts.\ 260--261 COIP; LM arts.\ 14, 17, 57; art.\ 316 CRE (delegación al sector privado implica distinción entre actividad legal e ilegal); Informe ARCOM 2024 &
La reforma mantiene con precisión la distinción entre minería privada formal, informal e ilegal; no extiende los beneficios de la apertura a operadores ilegales sin proceso de formalización; no crea ambigüedad que dificulte la persecución penal de los arts.\ 260--261 COIP; diferencia regímenes por escala &
La reforma define con precisión los sujetos que pueden acceder a concesiones. Existen regímenes diferenciados por escala. No hay vías de blanqueo de actividades ilegales. Existen sanciones proporcionales sin autorización \\
\hline

16 &
Incidencia en el sector minero privado &
Arts.\ 313, 316, 317, 339 CRE; LM arts.\ 1--2; Política Pública Minera Ecuador 2019; Plan Nacional Sector Minero 2020--2030 &
La reforma es compatible con la delegación excepcional del art.\ 316 CRE; contempla incentivos para la inversión formal sin reducir estándares ambientales como mecanismo de atracción; garantiza que los beneficios económicos para el Estado sean proporcionales a los impactos asumidos; no genera derechos adquiridos que impidan ajustes ambientales futuros ante nueva evidencia científica o jurisprudencial &
Los incentivos para inversión formal no están condicionados a la reducción de controles ambientales. Las regalías y participación del Estado son proporcionales a la escala. La reforma no extiende la delegación privada más allá del marco del art.\ 316 CRE \\
\hline

\end{longtable}
\endgroup

\noindent\textit{Nota.} Los dieciséis criterios derivan de normas constitucionales, legales y jurisprudenciales vigentes en el ordenamiento ecuatoriano, complementadas con estándares convencionales internacionales. Los criterios 6 (consulta ambiental, art.\ 398 CRE) y 7 (consulta previa, libre e informada, art.\ 57.7 CRE) son independientes entre sí: son mecanismos jurídicos distintos con fundamentos, procedimientos y sujetos diferentes, conforme a la Sentencia No.\ 22-18-IN/21 de la Corte Constitucional del Ecuador. El cumplimiento de uno no sustituye ni satisface el otro. El nivel de cumplimiento se determina como: \textit{cumple} / \textit{cumple parcialmente} / \textit{no cumple} / \textit{no aplica}. La Ley Orgánica para el Fortalecimiento de los Sectores Estratégicos de Minería y Energía (R.O.\ Quinto Suplemento No.\ 234, 2 de marzo de 2026) es empleada como referencia ilustrativa de aplicación de la Matriz; al momento de redacción de esta tesis registraba al menos once demandas de inconstitucionalidad ante la Corte Constitucional del Ecuador, cuyo resultado no ha sido determinado \parencite{primicias_demandas_2026}.
